\section{Introduction}
\label{sec:intro}




Efficiently and effectively endowing robotic policies with memory requires multiple levels of abstraction. While in principle we could simply encode the entire sequence of past observations into the context of the policy, this becomes intractable for long tasks, necessitating either very short sequences or significant subsampling.
However, in practical settings, the \emph{representation} required for long- and short-term memory is likely to be very different. For example, the robot might need to remember a recent observation to handle occlusions, and it might need to remember that it has already added one of the ingredients when cooking a meal. But these memories are fundamentally different: the former might require storing a few images over a short time period, the latter might require long-term memory but only a few bits of information.

An effective memory architecture for robot policies should use multiple modalities to represent memories at these different levels of abstraction. For short-horizon memory, dense image-based memory is well-suited to resolve occlusions and allows the robot to quickly adapt its manipulation strategy, e.g., by changing the grasp after failing to pick up an object. For long-horizon memory, we often only need to keep track of events at a \emph{semantic} level, such as which ingredient has already been added to a dish. In this case, a language-based representation provides much better compression than raw observations, and allows us to store high-level memories over long time periods.

Based on these observations, we introduce \MethodName{} (\Acronym{}), a system for equipping policies with multi-modal, long-horizon memory. \Acronym{} combines two key ingredients to make long-horizon memory tractable. 
First, we use a video encoder architecture to effectively encode multiple seconds of dense image-based memory into a compact representation. Second, we introduce a language-based memory mechanism in which the policy keeps track of semantic events in a compressed language format.
This memory system can not only accommodate very long horizon tasks, but also enables a variety of new capabilities by leveraging the short-term memory, such as in-context adaptation to correct mistakes, and resilience to partial observability and self-occlusion. 

To evaluate \Acronym{}, we integrate it into the \pizerosix{} model~\citep{pi06_model_card}, a generalist VLA trained on a diverse mixture of robot, vision-language, and video data.
We show that the resulting policy achieves state-of-the-art performance across a wide range of complex manipulation tasks. We also show that \Acronym{} enables our policy to solve long-horizon tasks like cleaning up a whole kitchen or preparing a grilled cheese sandwich, which require keeping track of memories for up to fifteen minutes.

In summary, we introduce a system for multi-scale, long-horizon memory for robot policies. By effectively representing short- and long-horizon memory via video and language representations respectively, \Acronym{} allows robot policies to keep track of memories across tens of minutes, without sacrificing runtime constraints. To implement \Acronym{}, we use an efficient video encoder architecture and a language-based memory system, and demonstrate their effectiveness through state-of-the-art performance across diverse robot tasks. With \Acronym{}, we enable robot policies to perform complex tasks like cleaning up whole kitchens, which can span up to fifteen minutes.























